\chapter{Foundations on unibranch compact complex spaces}
\label{app:hermitian-spaces}

A potential descends from a resolution when it is constant on each fiber. For a pulled-back background, local smooth correction and the maximum principle give this constancy on connected fibers. Descent then transports envelopes and compactness to the singular space, while vanishing on exceptional sets makes non-pluripolar measures independent of the resolution.

Throughout this appendix $X$ is a reduced irreducible compact unibranch complex space of pure dimension $n\geq1$. Smooth functions and forms on $X$ are understood through compatible local embeddings in complex Euclidean spaces. A Hermitian form is locally the restriction of a smooth positive ambient form. Fix a smooth real $(1,1)$-form $\theta$ on $X$.

Resolutions and common smooth refinements exist, and a proper bimeromorphic map to a unibranch space has connected fibers; see \cite[Theorems B.1.1--B.1.3]{SGPT}. The Fornaess--Narasimhan criterion identifies plurisubharmonic functions by upper semicontinuity and subharmonicity on holomorphic disks \cite[Theorem B.2.1]{SGPT}. The ramified disk lift used below is the argument in \cite[proof of Theorem B.2.2]{SGPT}.

\section{Potentials and disk tests}

\label{sec:space-background-setting}
The pullback of $\theta$ to the unit disk $\Delta$ has a smooth local potential. Disk tests therefore express the background inequality even when $\theta$ has no local potential on $X$.

\begin{definition}\label{def:space-theta-psh}
An upper semicontinuous function $u:X\to[-\infty,+\infty)$, not identically $-\infty$ on any nonempty open subset, is \emph{$\theta$-plurisubharmonic} if the following disk condition holds. For every holomorphic $f:\Delta\to X$, the function $u\circ f+a$ is subharmonic or identically $-\infty$ wherever $a$ is a smooth local potential for $f^*\theta$. The set of these functions is denoted by $\PSH(X,\theta)$.
\end{definition}

Two choices of $a$ differ by a harmonic function, so the condition is independent of that choice. On every disk where the pullback is genuine, it gives $\ddc(u\circ f)+f^*\theta\geq0$ in distributions.

\begin{lemma}\label{lem:space-local-correction}
Every member of $\PSH(X,\theta)$ is locally the sum of a plurisubharmonic function on $X$ and a smooth function. If $X$ is smooth, this definition agrees with \cref{def:theta-psh}.
\end{lemma}
\begin{proof}
In a local embedding choose an ambient smooth strictly plurisubharmonic function $r$ with $\ddc r\geq\theta$ on a smaller neighborhood. For a disk $f$ and a local potential $a$ of $f^*\theta$, the function $r\circ f-a$ is subharmonic. Hence $(u+r)\circ f$ is subharmonic or identically $-\infty$. The Fornaess--Narasimhan criterion and the local nondegeneracy of $u$ make $u+r$ plurisubharmonic. The smooth case now follows from \cref{lem:smooth-disk-criterion}.
\end{proof}

\section{Descent through a resolution}

The pullback of the background has no vertical component. A local correction therefore makes a potential plurisubharmonic near a fiber, and connectedness identifies the constants on its irreducible components.

\begin{theorem}\label{thm:hermitian-space-descent}
Let $\pi:Y\to X$ be a resolution. Pullback is an order-preserving bijection
\[
 \pi^*:\PSH(X,\theta)\longrightarrow\PSH(Y,\pi^*\theta).
\]
It preserves suprema of individual functions, boundedness, and ordinary singularity comparisons.
\end{theorem}
\begin{proof}
Pullback preserves the disk inequalities, and its genuineness follows from the isomorphism over a dense open subset. Injectivity follows from surjectivity of $\pi$.

For the converse take $v\in\PSH(Y,\pi^*\theta)$. Near a fiber choose the local correction $r$ on $X$ used above. The function $v+r\circ\pi$ is plurisubharmonic on its inverse image, since
\[
 \ddc(v+r\circ\pi)\geq\pi^*(\ddc r-\theta)\geq0.
\]
Its restriction to each compact irreducible component of the fiber is constant, possibly identically $-\infty$, by the maximum principle on a resolution of that component. Connectedness of the fiber makes the constants equal: intersecting components have the same value, and the finite component intersection graph is connected. Since $r\circ\pi$ is constant on the fiber, so is $v$.

Thus $v=u\circ\pi$ for a unique function $u$ on $X$. It is upper semicontinuous: for each real $c$, the closed set $\{v\geq c\}$ has closed image under the proper map, and that image is $\{u\geq c\}$. If $u$ were identically $-\infty$ on a nonempty open set, the same would hold for $v$ on its open inverse image. Thus $u$ is locally genuine.

It remains to retain the exact background in the disk test. Let $f:\Delta\to X$ have a genuine pullback $u\circ f$. Near a chosen point, take a curve in $\Delta\times_XY$ dominating the disk. Its normalization gives, after shrinking, a finite cover $\rho(t)=t^k$ and a holomorphic lift $f':\Delta\to Y$ with $\pi\circ f'=f\circ\rho$. Choose a smooth local function $a$ on the original disk with $\ddc a=f^*\theta$. The function
\[
 (u\circ f+a)\circ\rho=v\circ f'+a\circ\rho
\]
is subharmonic by \cref{lem:smooth-disk-criterion}. Subharmonicity descends through $\rho$ by the finite-disk argument following that lemma. Thus $u\circ f+a$ is subharmonic at every point, proving the bijection.

Surjectivity preserves suprema, uniform bounds and inequalities with constants, proving the remaining assertions.
\end{proof}

\section{Envelopes and normalized compactness}

Fix a resolution $\pi:Y\to X$. The envelope constructions of \cref{sec:obstacle-envelopes,def:singularity-envelope} now apply on $X$. In the following obstacle statement the inequality against $f$ is pointwise, and its family of minorants is assumed nonempty and locally uniformly bounded above. The descent bijection identifies the competitors; the remaining issue is whether upper semicontinuous regularization commutes with pullback.

\Needspace{10\baselineskip}
\begin{proposition}\label{prop:space-envelope-pullback}
For such an extended-real obstacle $f:X\to[-\infty,+\infty]$, the envelope is genuine and
\[
 P_{\pi^*\theta}(f\circ\pi)=P_\theta(f)\circ\pi.
\]
It also holds for the ordinary model envelope $P_\theta[u]$, for every $u\in\PSH(X,\theta)$.
\end{proposition}
\begin{proof}
The preceding bijection identifies the competitor families. Let $s$ be their pointwise supremum on $X$; its pullback is the pointwise supremum upstairs. Its regularization on $Y$ is a genuine $\pi^*\theta$-plurisubharmonic function, since there is a genuine lower competitor. It descends to a $\theta$-plurisubharmonic function $q$ by the theorem. Since $q\geq s$ and $q$ is upper semicontinuous, $q\geq s^*$.

For the reverse inequality fix $x\in X$ and $y\in\pi^{-1}(x)$. For each neighborhood $U$ of $x$, the unregularized supremum upstairs is bounded by $\sup_U s$ on the open neighborhood $\pi^{-1}(U)$ of $y$. Therefore $q(x)\leq\sup_U s$. Taking the infimum over $U$ gives $q(x)\leq s^*(x)$. This proves both the identity and genuineness downstairs. The proof for $P_\theta[u]$ is identical, because ordinary singularity order and its constants are preserved by pullback.
\end{proof}

The normalized space $\{u\in\PSH(X,\theta):\sup_Xu=0\}$ is compact in the $L^1$ topology of its pullbacks to $Y$, by \cref{thm:normalized-compactness} and descent. Smooth positive volume forms give equivalent norms on a fixed compact resolution. Under a further modification, pullback is continuous by \cref{lem:modification-lone-pullback} and bijective by descent. It is therefore a homeomorphism between the normalized compact spaces. A common resolution compares any two choices. This proves independence of the topology without asserting a comparison of the norms on different models.

\section{Ideal envelopes defined on resolutions}

Multiplier ideals can be compared on smooth models even when no intrinsic ideal sheaf on a singular local ring has been specified. Call two potentials on $X$ \emph{resolution-ideal equivalent} if their pullbacks are $\mathcal I$-equivalent on one smooth resolution. The all-weight equivalence of \cref{lem:ideal-equivalence-modification}, applied on a common refinement, makes this independent of the resolution. The smooth projection of \cref{thm:hermitian-i-projection} then defines an envelope in this class.

\Needspace{12\baselineskip}
\begin{proposition}\label{prop:space-resolution-i-envelope}
For $u\in\PSH(X,\theta)$, there is a unique potential $P_\theta[u]_{\mathcal I}^{\mathrm{res}}$ on $X$ satisfying
\[
 P_\theta[u]_{\mathcal I}^{\mathrm{res}}\circ\pi
       =P_{\pi^*\theta}[u\circ\pi]_{\mathcal I}
\]
on every resolution. It preserves the resolution-ideal class, is idempotent, has supremum zero, and is independent of the resolution.
\end{proposition}
\begin{proof}
For a fixed resolution the right side exists by \cref{thm:hermitian-i-projection} and descends by \cref{thm:hermitian-space-descent}. Under a further smooth modification, the ideal-envelope competitor families correspond by the descent theorem and \cref{lem:ideal-equivalence-modification}. The regularized suprema commute by the argument of \cref{prop:space-envelope-pullback}. Thus both resolutions give the same pullback on a common refinement and hence the same function on $X$. Projection, preservation of the class and normalization follow from the corresponding smooth assertions and surjectivity.
\end{proof}

\section{Extended measures and finite damping}

Write $X_{\mathrm{reg}}$ and $X_{\mathrm{sing}}$ for the regular and singular loci. Choose the resolution to be an isomorphism over $X_{\mathrm{reg}}$. The extended measure is determined on this isomorphism locus because it vanishes on the exceptional set. Its pushforward can thus be compared across resolutions even when its total mass is infinite.

For $u\in\PSH(X,\theta)$ define
\begin{equation}\label{eq:space-extended-ma}
 \mathrm{MA}^{\mathrm{ext}}_{X,\theta}(u)
   =\pi_*\mathrm{MA}^{\mathrm{ext}}_{Y,\pi^*\theta}(u\circ\pi),
\end{equation}
where $\mathrm{MA}^{\mathrm{ext}}_{Y,\pi^*\theta}$ is the extended measure of \cref{thm:arbitrary-extended-ma} on $Y$.

\enlargethispage{2\baselineskip}
\begin{theorem}\label{thm:space-extended-measures}
The right side is independent of the resolution. It is a positive sigma-finite Borel measure, gives zero mass to $X_{\mathrm{sing}}$ and to locally pluripolar sets, and may have infinite total mass. For bounded $u$ it is finite. Put $r_n=\max(n-2,0)$. There are finite constants $C$ and, for each $q>r_n$, $C_q$, depending only on the fixed background data and on $q$ in the second case, such that for every nonpositive $u\in\PSH(X,\theta)$ and $p>0$,
\[
 \int_X\mathrm{e}^{pu}\,\mathrm{MA}^{\mathrm{ext}}_{X,\theta}(u)
       \leq C(1+p^{-1})^{r_n},\qquad
 \int_X(1-u)^{-q}\,\mathrm{MA}^{\mathrm{ext}}_{X,\theta}(u)
       \leq C_q.
\]
The constants may be obtained on any one fixed resolution and are uniform over all such $u$.
\end{theorem}
\begin{proof}
On a resolution the preimage of $X_{\mathrm{sing}}$ is a proper analytic set. The smooth extended product gives it zero mass. Off that set, its construction is the intrinsic bounded local expansion on $X_{\mathrm{reg}}$, followed by truncation and exhaustion. Any two resolutions therefore give the same measure on the regular locus and zero measure on its complement. This proves independence.

Positivity follows upstairs. Sigma-finiteness also follows from the smooth construction: on the regular locus the resolution is a Borel isomorphism, so the finite-measure sets there push forward to a countable finite-measure exhaustion; the remaining analytic set has measure zero. Locally pluripolar sets pull back to pluripolar sets off the exceptional locus, and the exceptional locus itself has zero mass. Thus the stated vanishing holds.

If $u$ is bounded, its pullback is bounded on the compact smooth resolution. Local bounded-product estimates give finite total mass. For general nonpositive $u$, apply \cref{thm:general-background-damping} on $Y$ to the arbitrary smooth background $\pi^*\theta$. Pushforward and $u\circ\pi$ identify the weighted integrals exactly. The constants come from the fixed smooth background and Hermitian metric on $Y$.
\end{proof}

Define capacity convergence on $X$ by capacity convergence of pullbacks to a resolution, using a smooth Hermitian metric there. The local capacity comparison of \cref{sec:hermitian-capacity} makes the metric immaterial, and \cref{lem:modification-capacity-pullback} applied on a common resolution makes the choice of model immaterial.

It follows from \cref{thm:general-background-damping} and proper pushforward that both damped measures in the theorem are weakly continuous under this capacity convergence, for nonpositive potentials. For a sequence of nonpositive potentials $u_j\in\PSH(X,\theta)$ converging in capacity, the undamped measure is weakly continuous if, in addition, its total masses are uniformly bounded and its masses on $\{u_j<-T\}$ tend uniformly to zero as the real number $T$ tends to infinity. Pullback preserves these mass and sublevel conditions.

\section{Degenerate equations and singular traces}

Descent of potentials does not supply an equation estimate on the resolution. A pulled-back Hermitian form is generally degenerate there, and a pulled-back smooth volume can vanish along the exceptional locus. Existence, regularity and normalization therefore require estimates adapted to those degeneracies, as discussed in \cref{qu:singular-space-equations-program}.

An intrinsic restricted-ideal trace on a subspace contained in the singular locus requires regularization compatible with restriction there. The resolution-ideal envelope alone supplies neither this regularization nor the trace.
